\documentclass[tikz,border=10pt]{standalone}
\usepackage{fontspec}
\usepackage[mathrm=sym]{unicode-math}
\setmainfont{texgyrepagella-regular.otf}[
    BoldFont = texgyrepagella-bold.otf,
    ItalicFont = texgyrepagella-italic.otf,
    BoldItalicFont = texgyrepagella-bolditalic.otf
]
\setmathfont{texgyrepagella-math.otf}[AutoFakeBold = 2.5]
\usepackage{tikz}
\usepackage{pgfplots}
\pgfplotsset{compat=1.18}
\usetikzlibrary{arrows.meta, calc}

\begin{document}
\begin{tikzpicture}[>=Stealth, line cap=round, line join=round]

  % Scale factors: r in fm, V in MeV
  \def\xscale{0.36}
  \def\yscale{0.080}

  % Panel width
  \def\pw{11.5*\xscale}
  \def\gap{0.35}

  % Woods-Saxon potential curve definition (coordinates)
  \def\wsCurve{
    plot[smooth] coordinates {
      (0*\xscale, -44.0*\yscale)
      (3.0*\xscale, -44.0*\yscale)
      (4.5*\xscale, -43.8*\yscale)
      (5.5*\xscale, -42.2*\yscale)
      (6.5*\xscale, -36.5*\yscale)
      (7.2*\xscale, -27.5*\yscale)
      (7.8*\xscale, -17.5*\yscale)
      (8.4*\xscale, -8.5*\yscale)
      (9.0*\xscale, -3.8*\yscale)
      (9.8*\xscale, -1.0*\yscale)
      (10.5*\xscale, -0.2*\yscale)
    }
  }

  % --------------------------------------------------------
  % PANEL 1 (LEFT)
  % --------------------------------------------------------
  \begin{scope}[shift={(0, 0)}]

    % Top horizontal axis (V = 0)
    \draw[->, line width=0.85pt] (0, 0) -- (\pw, 0);

    % Left vertical axis (V(r))
    \draw[->, line width=0.85pt] (0, -45*\yscale) -- (0, 6*\yscale) node[above, font=\normalsize] {$V(r)$};

    % Y-axis ticks and labels
    \foreach \y in {-10, -20, -30, -40} {
      \draw[line width=0.65pt] (0, \y*\yscale) -- (0.15, \y*\yscale);
      \node[anchor=east, font=\small] at (-0.08, \y*\yscale) {$\y$};
    }

    % X-axis ticks (pointing up)
    \foreach \x in {5, 10} {
      \draw[line width=0.65pt] (\x*\xscale, 0) -- (\x*\xscale, 0.15);
      \node[anchor=south, font=\small] at (\x*\xscale, 0.10) {$\x$};
    }

    % Woods-Saxon potential (Solid)
    \draw[line width=1.15pt] \wsCurve;

    % Parabola tangent at r0 = 5.2 (Dotted)
    \draw[line width=1.2pt, dotted]
      plot[smooth, domain=0.2:8.8, samples=30]
      (\x*\xscale, {(-43.4 + 1.35 * (\x - 4.5)^2)*\yscale});

  \end{scope}

  % --------------------------------------------------------
  % PANEL 2 (MIDDLE)
  % --------------------------------------------------------
  \begin{scope}[shift={(\pw + \gap, 0)}]

    % Top horizontal axis (V = 0)
    \draw[->, line width=0.85pt] (0, 0) -- (\pw, 0);

    % Left border line
    \draw[line width=0.85pt] (0, -45*\yscale) -- (0, 0);

    % Y-axis ticks
    \foreach \y in {-10, -20, -30, -40} {
      \draw[line width=0.65pt] (0, \y*\yscale) -- (0.15, \y*\yscale);
    }

    % X-axis ticks (pointing up)
    \foreach \x in {5, 10} {
      \draw[line width=0.65pt] (\x*\xscale, 0) -- (\x*\xscale, 0.15);
      \node[anchor=south, font=\small] at (\x*\xscale, 0.10) {$\x$};
    }

    % Woods-Saxon potential (Solid)
    \draw[line width=1.15pt] \wsCurve;

    % Parabola tangent at r0 = 7.1 (Dotted) - restricted domain
    \draw[line width=1.2pt, dotted]
      plot[smooth, domain=2.6:8.25, samples=30]
      (\x*\xscale, {(-46.5 + 4.13 * (\x - 5.1)^2)*\yscale});

  \end{scope}

  % --------------------------------------------------------
  % PANEL 3 (RIGHT)
  % --------------------------------------------------------
  \begin{scope}[shift={(2*\pw + 2*\gap, 0)}]

    % Top horizontal axis (V = 0)
    \draw[->, line width=0.85pt] (0, 0) -- (\pw + 0.3, 0) node[right=2pt, font=\normalsize] {$r\,[\mathrm{fm}]$};

    % Left border line
    \draw[line width=0.85pt] (0, -45*\yscale) -- (0, 0);

    % Y-axis ticks
    \foreach \y in {-10, -20, -30, -40} {
      \draw[line width=0.65pt] (0, \y*\yscale) -- (0.15, \y*\yscale);
    }

    % X-axis ticks (pointing up)
    \foreach \x in {5, 10} {
      \draw[line width=0.65pt] (\x*\xscale, 0) -- (\x*\xscale, 0.15);
      \node[anchor=south, font=\small] at (\x*\xscale, 0.10) {$\x$};
    }

    % Woods-Saxon potential (Solid)
    \draw[line width=1.15pt] \wsCurve;

    % Inverted parabola tangent at r0 = 9.0 (Dotted)
    \draw[line width=1.2pt, dotted]
      plot[smooth, domain=6.5:10.9, samples=30]
      (\x*\xscale, {(-1.8 - 3.51 * (\x - 9.8)^2)*\yscale});

  \end{scope}

\end{tikzpicture}
\end{document}
